\documentclass{article}
\usepackage{amsmath}

\title{Q and P: Structured Comparator Drift and Dynamic Credit}
\author{}
\date{}

\begin{document}
\maketitle

\section*{The pair in two sentences}
Q studies LLM traders in a double auction and reports differences by model and
role, anchoring on visible offers, and an associational shift toward urgency when
one model closes the spread (ledger entries 2 and 7). P uses repeated multimodal
pairwise comparisons, Bradley--Terry rank aggregation, and an active-set
potential to assign credit in embodied multi-agent cooperation (entries 2, 3,
and 10).

\section*{The connexion pursued}
The proposed connexion is a robustness question: can P's repeated comparisons
remove noise when the comparator's preference changes systematically with order,
role, horizon cues, or coalition context? Q motivates those conditions as
diagnostics, but does not establish that they transfer from text-based traders
to multimodal referees (entries 7, 8, and 12). The final version separates
comparator drift from membership normalization because P shows visible energy
but does not clearly show horizon or the full active set to its comparator
(entries 18, 23, and 26).

\section*{What was found}
P's guarantee is narrower than convergence to true contribution. At a fixed
state it assumes comparisons are draws from one latent Bradley--Terry preference
and shows convergence to that LMM preference as the query count $K$ grows
(entries 2, 3, and 7). P reports pooled agreement with dense reward, while
prompt reversal addresses position bias; these checks do not establish
conditional stability across context interventions (entries 2 and 12).

The second round adds an input-level qualification. P's comparator receives two
ordered egocentric images and a task prompt; the images contain a visible energy
icon, and energy depletion governs removal and respawn. Exit proximity may
therefore be inferable, although remaining horizon and the full coalition are
not stated inputs (entries 23, 26, and 29). P also reports that increasing $K$
improves average downstream success. Hence the proposed study is not a generic
query-count ablation: it asks whether those pooled gains coexist with a
conditional bias floor (entry 32).

The notes formalize structured drift as
\begin{equation*}
 \Pr(Y_{ij}=1\mid z)=\frac{1}{1+\exp[-(c_i-c_j+d_{ij}(z))]}.
\end{equation*}
If a persistent distortion has the form $d_{ij}=b_i-b_j$, repeated queries
recover scores proportional to $c+b$: sampling variance can vanish while bias
remains. If $d_{ij}$ is not representable by score differences, cycles or
Bradley--Terry misfit can result (entry 11). Thus a larger $K$ may make credit
more repeatable without making it more valid; this is a derivation and testable
possibility, not an observed failure in P.

There is also a separate membership-boundary issue. P fixes its comparison
matrix at population size but defines scores and softmax over active agents
(entry 10). If survivor scores remain fixed, adding or removing a peer changes
the softmax denominator and therefore every survivor's normalized potential.
Successive potential differences can consequently contain a membership-boundary
impulse. The paper's identity alignment is underspecified, but the matrix
construction means a raw dimension error was not established (entries 10 and
14).

\section*{Objections and resolutions}
The peers rejected the broad claim that Q proves generic judge noise or causal
urgency: Q concerns autonomous traders, its urgency result is associational and
limited to one model, and its reasoning traces may be unfaithful (entries 2, 7,
and 12). They also corrected the claim that P necessarily has a dimension
mismatch and agreed that a well-defined Markov potential may retain standard
policy-invariance properties. The remaining concern is credit validity and
learning variance, not a theorem of policy bias (entries 10 and 14). Finally,
temporal variation alone does not contradict P because P permits a different
latent preference at each state; a valid test must hold task evidence fixed
while manipulating context within the comparisons for that state (entries 7
and 12).

The verifier accepted the connexion as promising but unverified and returned
ITERATE because no evidence transfers Q's trader behavior to P's multimodal
comparator (entry 22). The peers resolved how to phrase the energy observation:
masking or randomizing visible energy intervenes on an actual input, whereas
horizon or coalition text is added context. Neither may be called evidence of
membership-conditioned drift unless the frozen-state test shows it (entries 26,
29, and 30).

\section*{Outside sources}
No sources outside Q and P were used. The ledger and the two peer notes interpret only those papers.

\section*{What remains open}
The material is thin on outcomes. It is unknown whether P's code already aligns
fixed identities, whether membership information is present or inferable in the
LMM input, whether Q's context sensitivity generalizes to P's referee, whether
dense reward is an adequate contribution oracle, and whether any boundary
impulse changes learned policies (entries 16 and 18). Without a
context-conditioned bias floor, a downstream effect, and a remedy, the result
would be at most a robustness note (entries 11, 12, and 16).

The material remains thin: this round produced sharper controls but no test
result. Stage A is now a gate for the Q--P connexion; if it is null, Q supplies
no demonstrated mechanism. The normalization mechanism is independently gated
by a boundary impulse and downstream learning cost, and the mechanisms must not
be combined without evidence of interaction (entries 24, 27, and 30).

\section*{Smallest next step}
Freeze the pairwise task evidence at sampled P states and randomize only order,
identity, and an explicit horizon or coalition cue while crossing $K$. Measure
label flips, reversal consistency, cycles, held-out Bradley--Terry negative log
likelihood, and calibration to the chosen contribution reference. Separately
replay one matched removal or respawn transition through active-set and
entry/exit-invariant potentials and measure the boundary reward impulse
(entries 18 and 19). This small two-stage stress test settles whether either
proposed mechanism is present before a downstream training study is warranted.

The smallest immediate next step, as required by the verifier, is only stage A:
at one frozen P state, repeat comparisons across increasing $K$ while varying
image order, identity, and visible energy, with horizon or coalition cues kept
as separately labelled added-context conditions. A persistent conditional bias
floor would justify proceeding to the boundary replay; its absence would fail
the proposed Q--P connexion (entries 22, 26, 29, and 30).

\end{document}
